\documentclass{article}
\usepackage{amsmath}

\title{Pair Q7P7: Strategic Implementation of Forecast-Based Carbon Dispatch}
\author{Ada}
\date{}

\begin{document}
\maketitle

\section*{The pair in two sentences}
Q gives an indirect, budget-balanced mechanism and a proximal best-response
scheme for strongly concave welfare maximization over convex local strategy
sets with aggregate resource constraints.  P forecasts day-ahead nodal carbon
intensity and uses it to schedule distributed data centers (DDCs) and mobile
energy storage systems (MESSs), including binary intertemporal routing and
parking decisions, to reduce carbon emissions.

\section*{Connexion pursued}
The ledger asks whether Q can decentralize P's forecast-driven dispatch among
strategic asset owners and whether the resulting equilibrium inherits a
certificate against nodal-carbon-intensity forecast error.  The final form of
the question is deliberately narrower: embed forecasted carbon charges in Q's
mechanism on the largest compatible convex slice of P, rather than claim that
Q implements P's entire model (entries 2--4, 26, and 31).

\section*{What was found}
The full models do not compose directly.  Q assumes compact convex local sets,
twice differentiable strongly concave valuations, and shared constraints of
aggregate-capacity form; its implementation and uniqueness arguments use
these assumptions.  P instead minimizes a public carbon objective, and its
MESS model contains binary parking, travel, arrival, and departure variables.
Thus Q's KKT, LMI, uniqueness, and convergence conclusions do not presently
cover P's mixed-integer MESS dispatch (entries 2, 3, and 10).

There is nevertheless an algebraically valid mechanism on a compatible convex
slice.  Let owner \(n\)'s private benefit be \(v_n(x_n)\), let its separable
forecasted carbon contribution be
\(c_n(x_n)=\beta\widehat e^{\mathsf T}g_n(x_n)\), and let \(t_n^Q\) be Q's
transfer when Q is applied to the transformed valuation
\(V_n=v_n-c_n\).  Define
\[
 T_n=t_n^Q+c_n-\frac{1}{N-1}\sum_{m\ne n}c_m.
\]
Owner \(n\)'s payoff is
\[
 v_n-T_n=(v_n-c_n)-t_n^Q
 +\frac{1}{N-1}\sum_{m\ne n}c_m.
\]
The last term is independent of owner \(n\)'s message, so allocation and price
best responses are unchanged.  Each \(c_m\) appears in exactly \(N-1\)
rebates, hence
\[
 \sum_n(T_n-t_n^Q)=0.
\]
The augmentation therefore preserves Q's incentives and budget balance,
including off equilibrium, provided carbon is separable by owner allocation
(entries 24, 29, 32, and 38).  This construction is new to the pair; it is not
a theorem already stated by Q (entry 29).

This round resolves the verifier's principal objection for P's displayed
DDC-only feasible set.  At each time \(t\), write
\[
 w_{i,t}=\frac{\iota^{\mathrm{GLB}}_{i,t}}{\phi_{1i}},
 \qquad L_{i,t}\leq w_{i,t}\leq U_{i,t},
 \qquad \sum_i w_{i,t}=0,
\]
as in P's equations (27)--(30).  Introduce the nonnegative owner-local
coordinates
\[
 a_{i,t}=w_{i,t}-L_{i,t},\qquad
 b_{i,t}=U_{i,t}-w_{i,t},
\]
and put \(a_{i,t}+b_{i,t}=U_{i,t}-L_{i,t}\) in owner \(i\)'s local convex
set.  The two Q-form shared caps
\[
 \sum_i a_{i,t}\leq-\sum_iL_{i,t},\qquad
 \sum_i b_{i,t}\leq\sum_iU_{i,t}
\]
are respectively equivalent to \(\sum_iw_{i,t}\leq0\) and
\(\sum_iw_{i,t}\geq0\), and together are exactly P's signed conservation
equality.  Conversely, every P-feasible point saturates both caps.  This is
an affine bijection that preserves compactness and convexity, uses
nonnegative message coordinates, retains the DDC owners as agents, and
leaves P's DDC carbon term affine up to a baseline constant (entries 52, 56,
60, 61, and 69).  The capacities are nonnegative when baseline regulation is
feasible, so \(L_{i,t}\leq0\leq U_{i,t}\) (entries 61 and 69).  Strong
concavity in \(w_i\) also survives on the lifted local affine subspace, with
a rescaled modulus (entry 60).  Thus the connexion is demonstrated for P's
displayed DDC feasibility slice, assuming the private valuation required by
Q, rather than remaining wholly conditional (entries 63 and 72).

The forecast certificate is elementary but sharp within the stated model.
On a compact convex feasible set \(X\), define
\[
 F_e(x)=D(x)+\beta e^{\mathsf T}g(x),
\]
where \(D\) collects private costs.  Let the forecast and realized optimizers
be
\[
 \widehat x\in\operatorname*{argmin}_{x\in X}F_{\widehat e}(x),
 \qquad
 x_e^\star\in\operatorname*{argmin}_{x\in X}F_e(x).
\]
Forecast optimality and cancellation of the common private-cost terms give
\begin{align*}
 F_e(\widehat x)-F_e(x_e^\star)
 &\leq \beta(e-\widehat e)^{\mathsf T}
       \bigl(g(\widehat x)-g(x_e^\star)\bigr)\\
 &\leq \beta\lVert e-\widehat e\rVert_\infty
       \operatorname{diam}_1(g(X)).
\end{align*}
This improves the earlier \(2\beta\varepsilon B\) estimate when
\(\lVert g(x)\rVert_1\leq B\) (entries 5, 7, 9, 12, and 18).  If Q exactly
implements \(\widehat x\) on the compatible slice, the unique equilibrium has
the same realized-social-cost bound, so strategic decentralization contributes
no additional loss there (entries 18, 26, and 43).

\section*{Objections and their status}
Fixing MESS routes only linearizes products involving fixed binary parking
variables and does not settle any remaining storage complementarity or
additional grid coupling (entries 19, 21, and 70).  Entries 47, 53, and 55
initially treated P's signed DDC conservation as incompatible with Q, but the
paired-slack lift above resolves that objection: those peers retracted it
after noticing the essential owner-local complement equality (entries 60 and
61).  The lift covers P's displayed DDC equations (27)--(30), not any extra
shared grid polytope introduced in an implementation (entries 62 and 69).

Budget balance was initially said to prevent insertion of a public carbon
externality (entry 10).  The charge-and-rebate construction above resolves
that objection for separable carbon, and the other peers accepted the
correction (entries 24, 29, and 32).  Individual rationality is not similarly
settled.  Q uses a feasible zero action and the normalization \(V_n(0)=0\),
whereas P preserves mandatory baseline DDC service; zero regulation is
feasible but its incremental carbon may be signed.  MESS discharge credits
are signed as well.  Automatic inheritance of Q's participation guarantee is
therefore available only in an optional, nonnegative-consumption model with a
genuine zero-allocation outside option, not for P as written (entries 15, 30,
39, and 44).

The regret bound concerns realized social cost, not emissions alone, unless
service or utility is held fixed.  P's reported MAPE also does not directly
provide the required bus-time infinity-norm error, although its MaxAE is more
closely aligned (entries 8, 9, and 12).  Finally, Q's finite-iteration result
bounds fixed-point residuals but supplies neither a last-iterate rate nor an
error bound from residual to equilibrium distance.  An early-stopping term
needs an additional contraction or metric-subregularity argument (entry 42).

\section*{Sources outside Q and P}
No outside source was used.  The derivations above come only from Q, P, the
ledger, and the accompanying peer note.

\section*{What remains open}
It remains unknown whether P's DDC and grid constraints admit an exact
Q-compatible reformulation beyond the displayed DDC equations, whether
fixed-route MESS recourse retains a nonconvex charge/discharge condition, how
to prove participation relative to P's mandatory baseline and signed carbon
accounting, and how to turn Q's finite-iteration residual into operational
regret (entries 62, 67, 69, and 70).  The assumed smooth, strongly concave
private DDC valuation also lacks empirical or structural support in the
record (entries 64 and 70).  The material is thin on these points: there is no
participation proof, residual error bound, strategic calibration, or numerical
experiment.  Consequently the established contribution is still moderate
and domain-qualified, not a mechanism for P's full dispatch (entries 63, 67,
and 72).

\section*{Smallest settling step}
Prove individual rationality relative to P's baseline action in the lifted
DDC model.  The baseline \(w_{i,t}=0\) maps to
\(a_{i,t}=-L_{i,t}\), \(b_{i,t}=U_{i,t}\), while the zero vector used in Q's
participation deviation is infeasible unless \(U_{i,t}=L_{i,t}\).  A direct
baseline-relative proof, or incentive-neutral constants that remain budget
balanced, is now the smallest step that would settle whether the demonstrated
DDC implementation also has Q's participation property (entries 39, 62, 64,
and 72).

\end{document}
